\section{Examples and Definition}

\begin{definition}[Homomorphism]
If \(R\) and \(S\) are rings, a function \(f: R \rightarrow S\) is called a homomorphism if for all \(a, b \in R\) we have:
\begin{itemize}
\item[1.] \(f(1)=1\) (i.e., the identity in \(R\) is mapped to the one in \(S\); some people write this \(f\left(1_R\right)=1_S\).),
\item[2.] \(f(a+b)=f(a)+f(b)\),
\item[3.] \(f(a b)=f(a) f(b)\).
\end{itemize}

Just like for linear maps of vector spaces, homomorphisms have kernels and images. The kernel is

\[
\operatorname{Ker} f:=\{a \in R \mid f(a)=0\} \subseteq R,
\]

and the image is

\[
\operatorname{Im} f:=\{f(a) \mid a \in R\} \subseteq S .
\]
\end{definition}

\begin{lemma}
Let \(f: R \rightarrow S\) be a homomorphism of rings. Then
\begin{itemize}
\item[1.] \(f(0)=0\).
\item[2.] For all \(a \in R\) we have \(f(-a)=-f(a)\).
\end{itemize}
\end{lemma}

\begin{proofidea}
  
\end{proofidea}
  
\begin{formalproof}
  
\end{formalproof}

\begin{definition}[Isomorphic]
We say that two rings \(R\) and \(S\) are isomorphic if there exists a bijective homomorphism (called an isomorphism) between them and we then write \(R \cong S\).
\end{definition}

\begin{remark}
  To show that a map is an isomorphism, we have to show that it is a homomorphism, injective and surjective. Recall the definitions for injective and surjective:

\begin{itemize}
  \item \textbf{Injective} - \(f\left(x_1\right)=f\left(x_2\right) \Longrightarrow x_1=x_2\) 
  \item \textbf{Surjective} For every \(\mathrm{y} \in \mathrm{Y}\), there exists an \(\mathrm{x} \in \mathrm{X}\) such that \(\mathrm{f}(\mathrm{x})=\mathrm{y}\). This has \(Y\) be the cododmain.
\end{itemize}

\end{remark}

\begin{lemma}
A homomorphism \(f: R \rightarrow S\) is injective if and only if \(\operatorname{Ker} f=\{0\}\) (we will usually just write 0 for \(\{0\}\) ).
\end{lemma}

\begin{proofidea}
  
  
\end{proofidea}
  
\begin{formalproof}
  
\end{formalproof}

\section{Products of Rings and Projections}

\begin{definition}[Direct Product of Sets]
  
If we have two rings \(R\) and \(S\) we can construct a new ring, called the (direct) product of \(R\) and \(S\). It is denoted \(R \times S\) because as a set it is the Cartesian product of \(R\) and \(S\) (i.e., it consists of pairs \((r, s)\) where \(r \in R\) and \(s \in S\) ). Its binary operations are defined as follows:

\[
\begin{aligned}
(r, s)+\left(r^{\prime}, s^{\prime}\right) & :=\left(r+r^{\prime}, s+s^{\prime}\right), \\
(r, s)\left(r^{\prime}, s^{\prime}\right) & :=\left(r r^{\prime}, s s^{\prime}\right) .
\end{aligned}
\]

It is then clear that \((1,1) \in R \times S\) is the identity element. The other conditions for being a ring follow immediately from those for \(R\) and \(S\).

\end{definition}
